\documentclass[11pt]{article}
\usepackage[utf8]{inputenc}
\usepackage[T1]{fontenc}
\usepackage[italian]{babel}
\usepackage{geometry}
\geometry{a4paper, margin=1in}
\usepackage{xcolor}
\usepackage{listings}
\usepackage{hyperref}

\usepackage{tcolorbox}
\tcbuselibrary{breakable}
\begin{document}

\begin{tcolorbox}[colback=blue!5!white,colframe=blue!75!black,title=\textbf{DOMANDA}, breakable]
Esercizio
\begin{itemize}
    \item Considerare la cartella Webservice/
    \item Aprire il file \texttt{serverHTTP-REST.c} e analizzarne il contenuto
    \item Compilarlo come al solito ed eseguirlo
    \item Aprire il file \texttt{clientREST-GET.c} e analizzarne il contenuto
    \item Compilarlo come al solito ed eseguirlo
    \begin{itemize}
        \item Che parametri devo passare in linea di comando?
        \item Cosa si può vedere analizzando lo scambio di dati tramite Wireshark?
        \item Qual è la signature della funzione \texttt{calcolaSomma()} sul server e sul client? Perché ha senso che siano uguali?
    \end{itemize}
\end{itemize}
\end{tcolorbox}

\begin{tcolorbox}[colback=green!5!white,colframe=green!75!black,title=\textbf{RISPOSTA}, breakable]
\textbf{Soluzione Esercizio 29:} Webservice REST (RPC via HTTP)

\textbf{Operazioni Preliminari}\\
Ho copiato i file \texttt{serverHTTP-REST.c}, \texttt{clientREST-GET.c} e la libreria di rete all'interno della cartella dell'esercizio 29, compilandoli regolarmente con GCC.

\textbf{Risposte alle domande:}

\begin{enumerate}
    \item \textbf{Che parametri devo passare in linea di comando?}\\
    Guardando il codice del \texttt{main()} nel client, notiamo che l'eseguibile si aspetta che gli vengano passati 3 argomenti dopo il suo nome (\texttt{argc < 4}).\\
    \textbf{I parametri da passare sono, in ordine:} il nome della funzione da chiamare e i due operandi.\\
    \textbf{Esempio di utilizzo corretto:} \texttt{./clientREST-GET calcola-somma 10.5 4.5}

    \item \textbf{Cosa si può vedere analizzando lo scambio di dati tramite Wireshark?}\\
    Si osserva perfettamente il funzionamento di un'API REST:
    \begin{itemize}
        \item Il client effettua una richiesta \texttt{HTTP} GET formattando l'URL con i parametri "serializzati": \texttt{GET /calcola-somma?param1=10.500000\&param2=4.500000 HTTP/1.1}
        \item Il server risponde con un classico pacchetto \texttt{HTTP 200 OK}, ma invece di inviare codice HTML, nel Payload (Body) restituisce i dati strutturati formattati in \textbf{JSON}:
\begin{lstlisting}[language=json, basicstyle=\ttfamily\small]
{
    "somma": 15.000000
}
\end{lstlisting}
        Questo dimostra come i Webservice REST usino il protocollo web per lo scambio di dati "raw" tra macchine e non per le interfacce umane.
    \end{itemize}

    \item \textbf{Qual è la signature della funzione \texttt{calcolaSomma()} sul server e sul client? Perché ha senso che siano uguali?}\\
    \textbf{La signature (firma della funzione) è identica in entrambi i file:} \texttt{float calcolaSomma(float val1, float val2)}.\\
    Questo ha un grandissimo senso architetturale ed è alla base del concetto di \textbf{RPC (Remote Procedure Call)}:
    \begin{itemize}
        \item \textbf{Sul Server:} la funzione contiene la vera \textit{Business Logic}, ovvero svolge effettivamente il calcolo matematico \texttt{val1 + val2} e restituisce il risultato.
        \item \textbf{Sul Client:} la funzione funge da \textbf{Stub} (un intermediario/simulacro). Il suo scopo è "nascondere" al programmatore principale tutta la complessità della rete. Lo sviluppatore che scrive il \texttt{main()} del client può limitarsi a invocare \texttt{calcolaSomma(10, 5)} credendo di eseguire una normale funzione C locale. In realtà, lo Stub intercetta la chiamata, \textit{serializza} i parametri nell'URL, effettua la richiesta di rete \texttt{HTTP}, riceve la risposta, \textit{deserializza} il JSON per estrarne il numero puro, e lo restituisce al \texttt{main()}.
    \end{itemize}
\end{enumerate}

\textbf{Comandi utili:}
\begin{lstlisting}[language=bash, basicstyle=\ttfamily\small]
# Compila server e client
gcc serverHTTP-REST.c network.c -o serverHTTP-REST
gcc clientREST-GET.c network.c -o clientREST-GET

# Avvia il server su un terminale
./serverHTTP-REST

# Esegui il client su un altro terminale
./clientREST-GET calcola-somma 10.5 4.5
\end{lstlisting}
\end{tcolorbox}

\end{document}
